\section{Verification and open questions}\label{sec:verify}

Computer algebra enters in 3 places. In Lemma~\ref{lem:fisher-moments} each
moment lies in a space of sequences of dimension at most 10, and we check the
3 moments exactly for $N\le12$ with $q$ as a symbol, and the identity for
$\Var X^2$ in that lemma with sympy. The row $\Theta_\varepsilon$ of
Table~\ref{tab:outcomes} rests on 3 exact resultants of characteristic
polynomials. The values of $c_F$ on the 4-cycle after Corollary~\ref{cor:C4}
are exact evaluations of the next Laplace term of Remark~\ref{rem:nextterm}.
Every other number in the paper is a closed form or a computation at finite
size, done in 2 independent ways where possible. All 4 papers on Problem 131
share one verification supplement, \texttt{SUPPLEMENT.md} in the Problem 131
folder of \url{https://github.com/apemm/Kagey-Problems}. It describes the
checks, and \texttt{python verify\_all.py --paper c} in that folder runs the
ones for this paper. Before most computations we wrote a prediction
and a condition that would refute it in a ledger, and the outcome is recorded
below it. Many finite and algebraic steps, and the rows of
Table~\ref{tab:outcomes} marked Lean, are proved in Lean~4 in the folder
\texttt{lean/PaperC}, and its file \texttt{MAP.md} matches each statement of
the paper to its Lean theorems. The asymptotic results are not in Lean.

Several questions remain open. Theorem~\ref{thm:local} holds on compact
subsets of the open face, and a local law up to the boundary is open. The
same formula cannot hold there. Take a star with center $d$ and
$\ell=d-1\ge3$ leaves, and a start of full support. Then $P_N(n)=0$ whenever
$\supp n=[d]$ and $n_d<\ell-1$, since a path visits the center between any 2
runs at the leaves, although $n/N$ lies in the open face. This zero set is
checked in Lean. For a reducible face, Theorem~\ref{thm:first}(ii) places
the power of $T$ in $M_F$ between $k_F-k^*$ and $k_F-1$, and heuristically it
is $(k^*-1)/2+k_F-k^*$. Theorem~\ref{thm:crossing} puts all crossings of 2
faces in an interval of length $O(1/L)$, but uniqueness is proved only in
special cases, such as Papers~A and~B, the pairs of Proposition~\ref{prop:pair}
and the planar walk. We have no tree formula
for $K_F$ on non-reversible faces like Proposition~\ref{prop:constant}(b),
and no closed form for $c_F$. At finite $N$ the order of the phases can
differ from first order when a face lies close to the hull, as on the path
$P_4$ (Table~\ref{tab:outcomes}). For the planar walk the next term of the
crossing equation in Remark~\ref{rem:planar-turns} and the crossover near
$s=2r$ are open. A heuristic suggests that the next term comes from
replacing $T$ by $T+\frac1{4r}-\frac1{2(r+s)}$. For sticky priors we do not
treat random exit probabilities as in~\cite{fox2011}, or the probability of
a direct jump for $d\ge4$. Finally, we do not compute the constant in
Theorem~\ref{thm:fisher}, and the limits of $e_N$ at fixed $q$ and at fixed
$Nq$ are open.
